% ============================================================
%  ECON306: The Economics of Health and Education
%  University of Otago — Semester 2, 2026
%  Lecturer: Austin Henderson
%  LaTeX Beamer — Clean Minimal Academic Style
%  Compiled: \today
% ============================================================

\documentclass[aspectratio=169, 12pt]{beamer}

% ============================================================
%  ECON306: Shared Preamble — The Economics of Health and Education
%  University of Otago — Semester 2, 2026
%
%  SHARED BY ALL ECON306 DAY DECKS.
%  Input this file immediately after \documentclass in each deck:
%    % ============================================================
%  ECON306: Shared Preamble — The Economics of Health and Education
%  University of Otago — Semester 2, 2026
%
%  SHARED BY ALL ECON306 DAY DECKS.
%  Input this file immediately after \documentclass in each deck:
%    % ============================================================
%  ECON306: Shared Preamble — The Economics of Health and Education
%  University of Otago — Semester 2, 2026
%
%  SHARED BY ALL ECON306 DAY DECKS.
%  Input this file immediately after \documentclass in each deck:
%    \input{../Preambles/econ306_preamble.tex}
%  Do NOT include \documentclass or per-deck metadata here.
% ============================================================

% ─── Theme & Colours ────────────────────────────────────────
\usetheme{default}
\usecolortheme{default}
\usefonttheme{professionalfonts}
\definecolor{OtagoBlue}{HTML}{003057}
\definecolor{OtagoGold}{HTML}{A8923A}
\definecolor{TextGrey}{HTML}{3A3A3A}
\definecolor{AccentBlue}{HTML}{2B6CB0}
\definecolor{LightBg}{HTML}{F7F9FC}

\setbeamercolor{frametitle}{fg=OtagoBlue, bg=white}
\setbeamercolor{title}{fg=OtagoBlue}
\setbeamercolor{subtitle}{fg=OtagoGold}
\setbeamercolor{author}{fg=TextGrey}
\setbeamercolor{date}{fg=TextGrey}
\setbeamercolor{institute}{fg=TextGrey}
\setbeamercolor{normal text}{fg=TextGrey, bg=white}
\setbeamercolor{item}{fg=OtagoBlue}
\setbeamercolor{subitem}{fg=AccentBlue}
\setbeamercolor{block title}{bg=OtagoBlue!12, fg=OtagoBlue}
\setbeamercolor{block body}{bg=OtagoBlue!5, fg=TextGrey}
\setbeamercolor{alertblock title}{bg=OtagoGold!20, fg=OtagoBlue}
\setbeamercolor{alertblock body}{bg=OtagoGold!8, fg=TextGrey}
\setbeamercolor{alerted text}{fg=OtagoGold!80!black}
\setbeamercolor{structure}{fg=OtagoBlue}

\setbeamertemplate{frametitle}{%
	\vspace{0.35em}%
	\textbf{\insertframetitle}\par%
	\vspace{0.05em}%
	\textcolor{OtagoGold}{\rule{\textwidth}{0.6pt}}}

\setbeamertemplate{navigation symbols}{}
\setbeamertemplate{footline}{%
	\leavevmode\hbox{%
		\begin{beamercolorbox}[wd=.40\paperwidth,ht=2.5ex,dp=1.5ex,leftskip=.3cm]{author in head/foot}%
			\textcolor{TextGrey}{\footnotesize ECON306 \textbullet\ University of Otago}%
		\end{beamercolorbox}%
		\begin{beamercolorbox}[wd=.47\paperwidth,ht=2.5ex,dp=1.5ex,center]{title in head/foot}%
			\textcolor{TextGrey}{\footnotesize \insertshorttitle}%
		\end{beamercolorbox}%
		\begin{beamercolorbox}[wd=.13\paperwidth,ht=2.5ex,dp=1.5ex,rightskip=.3cm,right]{date in head/foot}%
			\textcolor{TextGrey}{\footnotesize \insertframenumber\,/\,\inserttotalframenumber}%
	\end{beamercolorbox}}\vskip0pt}

\setbeamertemplate{itemize item}{\textcolor{OtagoGold}{\small$\blacktriangleright$}}
\setbeamertemplate{itemize subitem}{\textcolor{AccentBlue}{\small$\triangleright$}}
\setbeamertemplate{enumerate item}{\textcolor{OtagoBlue}{\small\bfseries\insertenumlabel.}}

\setbeamertemplate{title page}{%
	\begin{tikzpicture}[remember picture,overlay]
		\fill[OtagoBlue] (current page.north west)
		rectangle ([yshift=-0.38\paperheight]current page.north east);
	\end{tikzpicture}
	\vspace{0.06\paperheight}
	\begin{beamercolorbox}[wd=\textwidth, sep=0pt]{title}
		{\color{white}\Large\textbf{\inserttitle}}\\[0.4em]
		{\color{OtagoGold}\large\insertsubtitle}
	\end{beamercolorbox}
	\vspace{1.4em}
	\begin{beamercolorbox}[wd=\textwidth, sep=0pt]{author}
		{\color{TextGrey}\normalsize\insertauthor}\\[0.2em]
		{\color{TextGrey}\small\insertinstitute}\\[0.2em]
		{\color{TextGrey}\small\insertdate}
\end{beamercolorbox}}

% ─── Packages ──────────────────────────────────────────────
\usepackage{tikz}
\usepackage{booktabs}
\usepackage{array}
\usepackage{graphicx}
\usepackage{hyperref}
\usepackage{amsmath}
\usepackage{amssymb}
\usepackage{colortbl}
\usepackage{multirow}
\usepackage{xcolor}
\hypersetup{colorlinks=true, linkcolor=AccentBlue, urlcolor=AccentBlue,
	citecolor=AccentBlue, pdfauthor={Austin Henderson},
	pdftitle={ECON306 S2 2026}}

% ─── Images folder ─────────────────────────────────────────
%  Place all slide images in the 'images/' subfolder next to each deck.
%  Usage: \includegraphics[width=0.85\textwidth]{figure_name}
\graphicspath{{images/}}

% ─── Reusable macros ───────────────────────────────────────
% Recap slide at start of each lecture
\newcommand{\recapslide}[1]{%
	\begin{frame}{Recap: What Did We Cover Last Time?}
		\begin{block}{Key points from last lecture}
			\begin{itemize}#1\end{itemize}
		\end{block}
\end{frame}}

% Plan slide at start of each lecture
\newcommand{\planslide}[2]{%
	\begin{frame}{Today's Plan — #1}
		\begin{columns}
			\column{0.6\textwidth}
			\begin{itemize}#2\end{itemize}
			\column{0.38\textwidth}
			\begin{alertblock}{Reading}
				\footnotesize See Brightspace (Aoroa) for required readings before this lecture.
			\end{alertblock}
		\end{columns}
\end{frame}}

% Summary slide at end of each lecture
\newcommand{\summaryside}[1]{%
	\begin{frame}{Lecture Summary}
		\begin{exampleblock}{Key takeaways}
			\begin{itemize}#1\end{itemize}
		\end{exampleblock}
		\vspace{0.5em}
		{\footnotesize\textcolor{OtagoGold}{\textbf{Brightspace (Aoroa):}} slides, readings, and any announcements posted there.}
\end{frame}}

% NZ spotlight box
\newcommand{\nzbox}[2]{%
	\begin{alertblock}{\includegraphics[height=0.8em]{nz_flag_icon}\hspace{0.3em} NZ Spotlight: #1}
		#2
\end{alertblock}}

% Tutorial callout
\newcommand{\tutorialbox}[1]{%
	\begin{block}{\textbf{Tutorial This Week}}
		#1
\end{block}}

% ─── Section dividers ──────────────────────────────────────
\AtBeginSection[]{
	\begin{frame}[plain]
		\begin{tikzpicture}[remember picture,overlay]
			\fill[OtagoBlue!95] (current page.north west)
			rectangle (current page.south east);
			\node[anchor=center, text=white, font=\Large\bfseries,
			text width=0.9\paperwidth, align=center]
			at (current page.center) {\insertsectionhead};
			\node[anchor=south, text=OtagoGold, font=\normalsize,
			text width=0.9\paperwidth, align=center]
			at ([yshift=2.5cm]current page.south) {ECON306 — University of Otago};
		\end{tikzpicture}
\end{frame}}

%  Do NOT include \documentclass or per-deck metadata here.
% ============================================================

% ─── Theme & Colours ────────────────────────────────────────
\usetheme{default}
\usecolortheme{default}
\usefonttheme{professionalfonts}
\definecolor{OtagoBlue}{HTML}{003057}
\definecolor{OtagoGold}{HTML}{A8923A}
\definecolor{TextGrey}{HTML}{3A3A3A}
\definecolor{AccentBlue}{HTML}{2B6CB0}
\definecolor{LightBg}{HTML}{F7F9FC}

\setbeamercolor{frametitle}{fg=OtagoBlue, bg=white}
\setbeamercolor{title}{fg=OtagoBlue}
\setbeamercolor{subtitle}{fg=OtagoGold}
\setbeamercolor{author}{fg=TextGrey}
\setbeamercolor{date}{fg=TextGrey}
\setbeamercolor{institute}{fg=TextGrey}
\setbeamercolor{normal text}{fg=TextGrey, bg=white}
\setbeamercolor{item}{fg=OtagoBlue}
\setbeamercolor{subitem}{fg=AccentBlue}
\setbeamercolor{block title}{bg=OtagoBlue!12, fg=OtagoBlue}
\setbeamercolor{block body}{bg=OtagoBlue!5, fg=TextGrey}
\setbeamercolor{alertblock title}{bg=OtagoGold!20, fg=OtagoBlue}
\setbeamercolor{alertblock body}{bg=OtagoGold!8, fg=TextGrey}
\setbeamercolor{alerted text}{fg=OtagoGold!80!black}
\setbeamercolor{structure}{fg=OtagoBlue}

\setbeamertemplate{frametitle}{%
	\vspace{0.35em}%
	\textbf{\insertframetitle}\par%
	\vspace{0.05em}%
	\textcolor{OtagoGold}{\rule{\textwidth}{0.6pt}}}

\setbeamertemplate{navigation symbols}{}
\setbeamertemplate{footline}{%
	\leavevmode\hbox{%
		\begin{beamercolorbox}[wd=.40\paperwidth,ht=2.5ex,dp=1.5ex,leftskip=.3cm]{author in head/foot}%
			\textcolor{TextGrey}{\footnotesize ECON306 \textbullet\ University of Otago}%
		\end{beamercolorbox}%
		\begin{beamercolorbox}[wd=.47\paperwidth,ht=2.5ex,dp=1.5ex,center]{title in head/foot}%
			\textcolor{TextGrey}{\footnotesize \insertshorttitle}%
		\end{beamercolorbox}%
		\begin{beamercolorbox}[wd=.13\paperwidth,ht=2.5ex,dp=1.5ex,rightskip=.3cm,right]{date in head/foot}%
			\textcolor{TextGrey}{\footnotesize \insertframenumber\,/\,\inserttotalframenumber}%
	\end{beamercolorbox}}\vskip0pt}

\setbeamertemplate{itemize item}{\textcolor{OtagoGold}{\small$\blacktriangleright$}}
\setbeamertemplate{itemize subitem}{\textcolor{AccentBlue}{\small$\triangleright$}}
\setbeamertemplate{enumerate item}{\textcolor{OtagoBlue}{\small\bfseries\insertenumlabel.}}

\setbeamertemplate{title page}{%
	\begin{tikzpicture}[remember picture,overlay]
		\fill[OtagoBlue] (current page.north west)
		rectangle ([yshift=-0.38\paperheight]current page.north east);
	\end{tikzpicture}
	\vspace{0.06\paperheight}
	\begin{beamercolorbox}[wd=\textwidth, sep=0pt]{title}
		{\color{white}\Large\textbf{\inserttitle}}\\[0.4em]
		{\color{OtagoGold}\large\insertsubtitle}
	\end{beamercolorbox}
	\vspace{1.4em}
	\begin{beamercolorbox}[wd=\textwidth, sep=0pt]{author}
		{\color{TextGrey}\normalsize\insertauthor}\\[0.2em]
		{\color{TextGrey}\small\insertinstitute}\\[0.2em]
		{\color{TextGrey}\small\insertdate}
\end{beamercolorbox}}

% ─── Packages ──────────────────────────────────────────────
\usepackage{tikz}
\usepackage{booktabs}
\usepackage{array}
\usepackage{graphicx}
\usepackage{hyperref}
\usepackage{amsmath}
\usepackage{amssymb}
\usepackage{colortbl}
\usepackage{multirow}
\usepackage{xcolor}
\hypersetup{colorlinks=true, linkcolor=AccentBlue, urlcolor=AccentBlue,
	citecolor=AccentBlue, pdfauthor={Austin Henderson},
	pdftitle={ECON306 S2 2026}}

% ─── Images folder ─────────────────────────────────────────
%  Place all slide images in the 'images/' subfolder next to each deck.
%  Usage: \includegraphics[width=0.85\textwidth]{figure_name}
\graphicspath{{images/}}

% ─── Reusable macros ───────────────────────────────────────
% Recap slide at start of each lecture
\newcommand{\recapslide}[1]{%
	\begin{frame}{Recap: What Did We Cover Last Time?}
		\begin{block}{Key points from last lecture}
			\begin{itemize}#1\end{itemize}
		\end{block}
\end{frame}}

% Plan slide at start of each lecture
\newcommand{\planslide}[2]{%
	\begin{frame}{Today's Plan — #1}
		\begin{columns}
			\column{0.6\textwidth}
			\begin{itemize}#2\end{itemize}
			\column{0.38\textwidth}
			\begin{alertblock}{Reading}
				\footnotesize See Brightspace (Aoroa) for required readings before this lecture.
			\end{alertblock}
		\end{columns}
\end{frame}}

% Summary slide at end of each lecture
\newcommand{\summaryside}[1]{%
	\begin{frame}{Lecture Summary}
		\begin{exampleblock}{Key takeaways}
			\begin{itemize}#1\end{itemize}
		\end{exampleblock}
		\vspace{0.5em}
		{\footnotesize\textcolor{OtagoGold}{\textbf{Brightspace (Aoroa):}} slides, readings, and any announcements posted there.}
\end{frame}}

% NZ spotlight box
\newcommand{\nzbox}[2]{%
	\begin{alertblock}{\includegraphics[height=0.8em]{nz_flag_icon}\hspace{0.3em} NZ Spotlight: #1}
		#2
\end{alertblock}}

% Tutorial callout
\newcommand{\tutorialbox}[1]{%
	\begin{block}{\textbf{Tutorial This Week}}
		#1
\end{block}}

% ─── Section dividers ──────────────────────────────────────
\AtBeginSection[]{
	\begin{frame}[plain]
		\begin{tikzpicture}[remember picture,overlay]
			\fill[OtagoBlue!95] (current page.north west)
			rectangle (current page.south east);
			\node[anchor=center, text=white, font=\Large\bfseries,
			text width=0.9\paperwidth, align=center]
			at (current page.center) {\insertsectionhead};
			\node[anchor=south, text=OtagoGold, font=\normalsize,
			text width=0.9\paperwidth, align=center]
			at ([yshift=2.5cm]current page.south) {ECON306 — University of Otago};
		\end{tikzpicture}
\end{frame}}

%  Do NOT include \documentclass or per-deck metadata here.
% ============================================================

% ─── Theme & Colours ────────────────────────────────────────
\usetheme{default}
\usecolortheme{default}
\usefonttheme{professionalfonts}
\definecolor{OtagoBlue}{HTML}{003057}
\definecolor{OtagoGold}{HTML}{A8923A}
\definecolor{TextGrey}{HTML}{3A3A3A}
\definecolor{AccentBlue}{HTML}{2B6CB0}
\definecolor{LightBg}{HTML}{F7F9FC}

\setbeamercolor{frametitle}{fg=OtagoBlue, bg=white}
\setbeamercolor{title}{fg=OtagoBlue}
\setbeamercolor{subtitle}{fg=OtagoGold}
\setbeamercolor{author}{fg=TextGrey}
\setbeamercolor{date}{fg=TextGrey}
\setbeamercolor{institute}{fg=TextGrey}
\setbeamercolor{normal text}{fg=TextGrey, bg=white}
\setbeamercolor{item}{fg=OtagoBlue}
\setbeamercolor{subitem}{fg=AccentBlue}
\setbeamercolor{block title}{bg=OtagoBlue!12, fg=OtagoBlue}
\setbeamercolor{block body}{bg=OtagoBlue!5, fg=TextGrey}
\setbeamercolor{alertblock title}{bg=OtagoGold!20, fg=OtagoBlue}
\setbeamercolor{alertblock body}{bg=OtagoGold!8, fg=TextGrey}
\setbeamercolor{alerted text}{fg=OtagoGold!80!black}
\setbeamercolor{structure}{fg=OtagoBlue}

\setbeamertemplate{frametitle}{%
	\vspace{0.35em}%
	\textbf{\insertframetitle}\par%
	\vspace{0.05em}%
	\textcolor{OtagoGold}{\rule{\textwidth}{0.6pt}}}

\setbeamertemplate{navigation symbols}{}
\setbeamertemplate{footline}{%
	\leavevmode\hbox{%
		\begin{beamercolorbox}[wd=.40\paperwidth,ht=2.5ex,dp=1.5ex,leftskip=.3cm]{author in head/foot}%
			\textcolor{TextGrey}{\footnotesize ECON306 \textbullet\ University of Otago}%
		\end{beamercolorbox}%
		\begin{beamercolorbox}[wd=.47\paperwidth,ht=2.5ex,dp=1.5ex,center]{title in head/foot}%
			\textcolor{TextGrey}{\footnotesize \insertshorttitle}%
		\end{beamercolorbox}%
		\begin{beamercolorbox}[wd=.13\paperwidth,ht=2.5ex,dp=1.5ex,rightskip=.3cm,right]{date in head/foot}%
			\textcolor{TextGrey}{\footnotesize \insertframenumber\,/\,\inserttotalframenumber}%
	\end{beamercolorbox}}\vskip0pt}

\setbeamertemplate{itemize item}{\textcolor{OtagoGold}{\small$\blacktriangleright$}}
\setbeamertemplate{itemize subitem}{\textcolor{AccentBlue}{\small$\triangleright$}}
\setbeamertemplate{enumerate item}{\textcolor{OtagoBlue}{\small\bfseries\insertenumlabel.}}

\setbeamertemplate{title page}{%
	\begin{tikzpicture}[remember picture,overlay]
		\fill[OtagoBlue] (current page.north west)
		rectangle ([yshift=-0.38\paperheight]current page.north east);
	\end{tikzpicture}
	\vspace{0.06\paperheight}
	\begin{beamercolorbox}[wd=\textwidth, sep=0pt]{title}
		{\color{white}\Large\textbf{\inserttitle}}\\[0.4em]
		{\color{OtagoGold}\large\insertsubtitle}
	\end{beamercolorbox}
	\vspace{1.4em}
	\begin{beamercolorbox}[wd=\textwidth, sep=0pt]{author}
		{\color{TextGrey}\normalsize\insertauthor}\\[0.2em]
		{\color{TextGrey}\small\insertinstitute}\\[0.2em]
		{\color{TextGrey}\small\insertdate}
\end{beamercolorbox}}

% ─── Packages ──────────────────────────────────────────────
\usepackage{tikz}
\usepackage{booktabs}
\usepackage{array}
\usepackage{graphicx}
\usepackage{hyperref}
\usepackage{amsmath}
\usepackage{amssymb}
\usepackage{colortbl}
\usepackage{multirow}
\usepackage{xcolor}
\hypersetup{colorlinks=true, linkcolor=AccentBlue, urlcolor=AccentBlue,
	citecolor=AccentBlue, pdfauthor={Austin Henderson},
	pdftitle={ECON306 S2 2026}}

% ─── Images folder ─────────────────────────────────────────
%  Place all slide images in the 'images/' subfolder next to each deck.
%  Usage: \includegraphics[width=0.85\textwidth]{figure_name}
\graphicspath{{images/}}

% ─── Reusable macros ───────────────────────────────────────
% Recap slide at start of each lecture
\newcommand{\recapslide}[1]{%
	\begin{frame}{Recap: What Did We Cover Last Time?}
		\begin{block}{Key points from last lecture}
			\begin{itemize}#1\end{itemize}
		\end{block}
\end{frame}}

% Plan slide at start of each lecture
\newcommand{\planslide}[2]{%
	\begin{frame}{Today's Plan — #1}
		\begin{columns}
			\column{0.6\textwidth}
			\begin{itemize}#2\end{itemize}
			\column{0.38\textwidth}
			\begin{alertblock}{Reading}
				\footnotesize See Brightspace (Aoroa) for required readings before this lecture.
			\end{alertblock}
		\end{columns}
\end{frame}}

% Summary slide at end of each lecture
\newcommand{\summaryside}[1]{%
	\begin{frame}{Lecture Summary}
		\begin{exampleblock}{Key takeaways}
			\begin{itemize}#1\end{itemize}
		\end{exampleblock}
		\vspace{0.5em}
		{\footnotesize\textcolor{OtagoGold}{\textbf{Brightspace (Aoroa):}} slides, readings, and any announcements posted there.}
\end{frame}}

% NZ spotlight box
\newcommand{\nzbox}[2]{%
	\begin{alertblock}{\includegraphics[height=0.8em]{nz_flag_icon}\hspace{0.3em} NZ Spotlight: #1}
		#2
\end{alertblock}}

% Tutorial callout
\newcommand{\tutorialbox}[1]{%
	\begin{block}{\textbf{Tutorial This Week}}
		#1
\end{block}}

% ─── Section dividers ──────────────────────────────────────
\AtBeginSection[]{
	\begin{frame}[plain]
		\begin{tikzpicture}[remember picture,overlay]
			\fill[OtagoBlue!95] (current page.north west)
			rectangle (current page.south east);
			\node[anchor=center, text=white, font=\Large\bfseries,
			text width=0.9\paperwidth, align=center]
			at (current page.center) {\insertsectionhead};
			\node[anchor=south, text=OtagoGold, font=\normalsize,
			text width=0.9\paperwidth, align=center]
			at ([yshift=2.5cm]current page.south) {ECON306 — University of Otago};
		\end{tikzpicture}
\end{frame}}


% ─── Document metadata ─────────────────────────────────────
\title[ECON306: Health \& Education Economics]{ECON306\\The Economics of Health and Education}
\subtitle{Semester 2, 2026}
\author{Austin Henderson}
\institute{Department of Economics\\University of Otago}
\date{Semester 2 — July to October 2026}

% ============================================================

\begin{document}

% ─── Title slide ───────────────────────────────────────────
\maketitle

% ================================================================
%  LECTURE 1 — Tuesday 14 July
% ================================================================

\begin{frame}{Week 1 Overview}
	\begin{columns}
		\column{0.55\textwidth}
		\textbf{Three lectures this week:}
		\begin{enumerate}
			\item What is health economics? Special characteristics of health care
			\item Positive vs.\ normative analysis; the HC ``problem'' in NZ
			\item Determinants of health; the health production function
		\end{enumerate}
		\column{0.42\textwidth}
		\begin{block}{Required Reading}
			{\footnotesize
				\begin{itemize}\setlength{\itemsep}{0.1em}
					\item Santerre \& Neun: ``Theories X and Y of health economics''
					\item Alchian \& Allen: ``Need versus demand''
					\item \textit{The Economist}, ``Catching Up'' (Jan 2013)
					\item Folland, Goodman \& Stano: Ch.\ 5
					\item Hansen \& Graham: ``Human organ transplants''
			\end{itemize}}
		\end{block}
	\end{columns}
\end{frame}

% ─── Section: What is Health Economics? ────────────────────
\section{What is Health Economics?}

\planslide{Lecture 1}{
	\item What is health economics? Definitions and scope
	\item Is health care an \emph{economic good}?
	\item The \textbf{six special demand \& supply characteristics} of health care
	\item Implications for how health care is funded and organised
	\item Introduction to the course approach: positive vs.\ normative
}

\begin{frame}{What is ``Health Economics''?}
	\begin{columns}
		\column{0.52\textwidth}
		\begin{block}{Three definitions}
			{\small\begin{itemize}\setlength{\itemsep}{0.1em}
				\item \textit{Kobelt (1996):} ``the application of theories, tools and concepts of economics to health and health care''
				\item \textit{NICE (UK):} ``the study of the cost of using and distributing health care resources''
				\item \textit{Wikipedia:} ``concerned with efficiency, effectiveness, value and behavior in the production and consumption of health and health care''
			\end{itemize}}
		\end{block}
		\column{0.44\textwidth}
		\begin{alertblock}{Economics, in brief}
			{\small The study of how individuals and societies use \textbf{scarce resources}.\\[0.3em]
			Faced with scarcity, how do we make the \textbf{best} choices?}
		\end{alertblock}
		\vspace{0.3em}
		\begin{exampleblock}{Why it matters for NZ}
			{\small Health + education $\approx$ \textbf{13\% of NZ's GDP}\\
			Health alone: $\approx$\textbf{\$30.6B}/year in 2024/25 ($\sim$8\% GDP)}
		\end{exampleblock}
	\end{columns}
\end{frame}

\begin{frame}{Four Fundamental Questions}
	\begin{columns}
		\column{0.62\textwidth}
		\begin{enumerate}
			\item Is health care an \textbf{economic good}?\\
			{\small \textcolor{AccentBlue}{Is health care a fundamental human right?}}
			\vspace{0.5em}
			\item How does health care \textbf{differ} from other economic goods?\\
			{\small \textcolor{AccentBlue}{Is health care ``special''?}}
			\vspace{0.5em}
			\item Can we study health and health care using \textbf{economic analysis}?\\
			{\small \textcolor{AccentBlue}{Supply and demand?}}
			\vspace{0.5em}
			\item What are the \textbf{implications} of these special characteristics for how health
			care is funded and organised?
		\end{enumerate}
		\column{0.35\textwidth}
		\begin{alertblock}{A running theme}
			These questions will recur in \textbf{every week} of the course.
		\end{alertblock}
	\end{columns}
\end{frame}

% ─── Section: The Six Special Characteristics ──────────────
\section{The Six Special Characteristics}

\begin{frame}{Special Characteristic 1: Uncertainty}
	\begin{columns}
		\column{0.58\textwidth}
		{\small\begin{itemize}\setlength{\itemsep}{0.2em}
			\item Classic economics assumes \emph{perfect information} --- but health care violates this
			\item \textbf{Demand-side}: How much HC will you need? When? What type?
			Unlike demand for laptops or t-shirts, you cannot know in advance.
			\item \textbf{Supply-side}: Doctors decide under uncertainty --- right diagnosis? Right treatment? What is best for a patient vs.\ a population with scarce resources?
		\end{itemize}}
		\column{0.38\textwidth}
		\begin{block}{Institutional response}
			{\small\textbf{Insurance} (public or private) pools risk across many people.}
		\end{block}
		\vspace{0.3em}
		\begin{alertblock}{But insurance creates problems}
			{\small\textit{Moral hazard} and \textit{adverse selection} (Week 5)}
		\end{alertblock}
	\end{columns}
\end{frame}

\begin{frame}{Special Characteristic 2: Derived Demand}
	\begin{columns}
		\column{0.55\textwidth}
		\begin{itemize}
			\item People don't want health \emph{care} for its own sake --- they want \textbf{health}
			\item HC is therefore a \textit{derived demand}: it is demanded only because it (sometimes)
			produces health
			\item Policy implication: if other inputs (e.g.\ diet, exercise, clean air) produce health
			more cost-effectively than HC, \emph{that is where resources should go}
			\item Demand signals may be weaker when demand is derived --- consumer behaviour is more
			passive
		\end{itemize}
		\column{0.42\textwidth}
		\begin{block}{Health Production Function}
			{\small$$H = f\bigl(\underbrace{\text{HC},\; \text{nutrition},\; \text{housing},\; \ldots}_{\text{inputs}}\bigr)$$}
			\vspace{0.2em}
			{\small Changing relative prices of any input affects the optimal mix.}
		\end{block}
	\end{columns}
\end{frame}

\begin{frame}{Special Characteristic 3: Asymmetric Information}
	\begin{columns}
		\column{0.55\textwidth}
		{\small\begin{itemize}\setlength{\itemsep}{0.2em}
			\item \textbf{Principal-agent problem}: providers know more than patients --- but are their incentives aligned?
			\item \textbf{Supplier-Induced Demand (SID)}: doctors may induce demand for their own services (profit motive)
			\item On the insurance side: high-risk patients seek comprehensive cover --- \textit{adverse selection}
			\item Societal responses: licensing, credentialing, standardised treatment protocols, transparency laws
		\end{itemize}}
		\column{0.42\textwidth}
		\begin{exampleblock}{NZ Spotlight}
			{\small NZ and the US are the only developed economies that allow \textbf{direct-to-consumer (DTC) pharmaceutical advertising}.\\[0.3em]
			Implications for SID --- see Week 3.}
		\end{exampleblock}
	\end{columns}
\end{frame}

\begin{frame}{Special Characteristic 4: Catastrophic \& Irreversible Consequences}
	\begin{columns}
		\column{0.55\textwidth}
		\begin{itemize}
			\item Health outcomes can be \textbf{life-threatening} at the extremes of the distribution
			\item Many HC decisions are \textbf{irreversible} --- health is very path-dependent
			\item Risk of catastrophe may matter more to individuals than expected-value averages
		\end{itemize}
		\column{0.42\textwidth}
		\begin{alertblock}{Why this matters}
			Unlike buying the wrong laptop, the wrong medical decision can have \textbf{permanent consequences}.\\[0.4em]
			This asymmetry in stakes changes how rational agents should behave.
		\end{alertblock}
	\end{columns}
\end{frame}

\begin{frame}{Special Characteristic 5: Externalities}
	\begin{columns}
		\column{0.52\textwidth}
		\begin{block}{Physical externalities}
			\begin{itemize}
				\item Immunisation: my vaccination protects you
				\item Communicable diseases: my illness endangers you
				\item Food safety and sanitation: collective goods
			\end{itemize}
		\end{block}
		\column{0.44\textwidth}
		\begin{block}{Psychic (caring) externalities}
			\begin{itemize}
				\item We care about the health of others (altruism, social norms)
				\item Willingness to pay for others' HC even when not directly affected
			\end{itemize}
		\end{block}
	\end{columns}
	\vspace{0.5em}
	\begin{alertblock}{Implication}
		Externalities mean private market outcomes \textbf{undervalue} health care ---
		the social benefit exceeds the private benefit.
	\end{alertblock}
\end{frame}

\begin{frame}{Special Characteristic 6: Price Discrimination}
	\begin{columns}
		\column{0.55\textwidth}
		\begin{block}{Three degrees}
			\begin{enumerate}
				\item \textbf{First-degree} (perfect): charge each consumer their maximum willingness to pay
				\item \textbf{Second-degree} (bulk/bundles): quantity discounts, insurance tiers
				\item \textbf{Third-degree} (by group): different prices for students, elderly, insured vs.\ uninsured
			\end{enumerate}
		\end{block}
		\column{0.42\textwidth}
		\begin{alertblock}{In practice}
			Insurers use third-degree price discrimination extensively.\\[0.4em]
			Regulation governs this in most countries --- but the extent varies dramatically.
		\end{alertblock}
	\end{columns}
\end{frame}

% ─── Section: Implications ─────────────────────────────────
\section{Implications for Government Involvement}

\begin{frame}{The Case for Government Involvement}
	\begin{columns}[T]
		\column{0.55\textwidth}
		{\small These six characteristics create \textbf{market failures} --- unregulated markets will not produce efficient or equitable outcomes.\\[0.3em]
		Three forms of government involvement:
		\begin{enumerate}\setlength{\itemsep}{0.1em}
			\item \textbf{Funding} (tax-financed public insurance, subsidies)
			\item \textbf{Provision} (public hospitals, publicly employed doctors)
			\item \textbf{Regulation} (licensing, drug pricing)
		\end{enumerate}
		\vspace{0.1em}
		In NZ: government funds $\approx$77\% of total health spending.}
		\column{0.42\textwidth}
		\begin{alertblock}{But government can also fail}
			{\small Market failure $\Rightarrow$ intervention \emph{may} improve outcomes, but can also make things worse.\\[0.2em]
			We will critically examine both sides.}
		\end{alertblock}
	\end{columns}
\end{frame}

% ─── Summary ───────────────────────────────────────────────

\summaryside{
	\item Health economics applies economic theory to the production and consumption of health and health care
	\item Health care has \textbf{six special characteristics} that differentiate it from other goods: uncertainty, derived demand, asymmetric information, catastrophic consequences, externalities, and price discrimination
	\item These characteristics create \textbf{market failures} that justify government involvement in financing, provision, and regulation
	\item In NZ, health care accounts for $\approx$8\% of GDP and 21\% of government spending
}

\begin{frame}{What's Next?}
	\begin{columns}
		\column{0.55\textwidth}
		\begin{block}{Lecture 2}
			{\small\begin{itemize}\setlength{\itemsep}{0.1em}
				\item Positive vs.\ normative analysis
				\item The health care ``problem'' in NZ
				\item How do we evaluate policy interventions?
			\end{itemize}}
		\end{block}
		\vspace{0.3em}
		\begin{block}{Lecture 3}
			{\small\begin{itemize}\setlength{\itemsep}{0.1em}
				\item Determinants of health
				\item The health production function in detail
			\end{itemize}}
		\end{block}
		\column{0.42\textwidth}
		\begin{alertblock}{Before next time}
			{\footnotesize\begin{itemize}\setlength{\itemsep}{0.1em}
				\item Read: Alchian \& Allen, ``Need versus demand''
				\item Read: \textit{The Economist}, ``Catching Up''
				\item Think about: Is health care a \textit{right} or an \textit{economic good}?
			\end{itemize}}
		\end{alertblock}
	\end{columns}
\end{frame}

\end{document}
